\section{Preliminaries}\label{sec:preliminaries}

We introduce the \rankop/\selectop primitives on \bvs and common succinct
tree-encoding schemes, review four state-of-the-art succinct tries this paper
builds upon, and summarize string compression schemes for the tail container.

\subsection{Rank/Select Primitives}

\Bvs are a fundamental building block of many succinct data structures. Succinct
tree navigation relies on the \rankop and \selectop primitives of \bvs, which
are defined as follows:
\begin{itemize}
  \item \proc{rank$_p$(i)}: return the number of pattern $p$'s in the first $i$ bits of the \bv.
  \item \proc{select$_p$(i)}: return the end of the $i$-th occurrence of pattern $p$ in the \bv.
\end{itemize}
Different \bvs support different patterns for \rankop and \selectop. Common
choices of $p$ include $1$, $0$ and $00$. For example, \proc{rank$_1$(i)}
returns the number of occurrences of ``1'' in \rev{positions} [0, i).

%Without any auxiliary data structures, the worst-case time complexities of the \rankop and
% \selectop operations are both $O(n)$, where $n$ is the size of the bitvector.

Index structures for \rankop and \selectop partition the \bv into fixed-sized
(e.g., 512-bit) \rev{\emph{basic blocks}} and store partial results so far for
each block from left-to-right~\cite{rank_select_1}. \rev{They also often include
  \emph{secondary indices} on smaller fixed-size blocks within each basic block
  for higher granularity and better performance.} These auxiliary structures
reduce \newmymarginpar{R2 O3/D6}\newrev{the worst-case cost} of \rankop and
\selectop to $O(1)$ with minimal (typically \rev{about} $5\%$) space overhead
\cite{surf,rank_select_0,rank_select_1,rank_select_2,rank_select_3,sds_overview,ds2i,optimized_rank,terark_bv}.

For the \rankop operation, the \defn{rank index} stores \mymarginpar{R2
  D14}\rev{\defn{samples}, or exact accumulative ranks before each block (e.g.,
  the number of 1s preceding each basic block), which enable fast exact
  responses to \rankop queries.} Answering a \rev{\proc{rank$_p$(i)} query}
involves combining the accumulative rank before the target block with the
(computed) rank in the remainder up to position \proc{i}.

\mymarginpar{Meta 1, R1 D1}\new{For example,~\figref{traditional-bv} illustrates
  a worked example of a \\ \proc{rank$_1$(12)} query on a bitvector where the rank
  index is built with block size 8. First, we look up the result of
  \proc{rank$_1$(8) = 4} in $O(1)$ time in the rank index because position 8 is
  the start of the basic block containing 12. Second, we calculate the remaining
  instances of the pattern in the range 8 - 12 in the bitvector. The combined
  result of the first and second step is 4 + 1 = 5, which is
  \proc{rank$_1$(12)}.}

\rev{The \selectop operation uses a \defn{select index} similar
  to the aforementioned rank index, but \selectop is more challenging than
  \rankop because a query may not know exactly which block is needed based on
  the input. For example, with 8-bit basic blocks, bit 12 is always in the
  second basic block, but the 12-th \emph{one bit} might be in the 100-th basic
  block!

  To support the \selectop operation, the \selectop index precomputes and stores
  samples,\mymarginpar{R2 D14} or exact results of \selectop queries at regular
  intervals (e.g., \proc{select$_p$(0)}, \proc{select$_p$(8)},
  \proc{select$_p$(16)}, \ldots). Answering a \selectop query involves finding
  the closest index entry to the left of the desired select query, then scanning
  sequentially using that as a starting point.}  \mymarginpar{Meta 1, R1
  D1}\new{For example,~\figref{traditional-bv} illustrates a worked example of a
  \proc{select$_1$(14)} query on a bitvector where the select index is built
  with block size 8. First, we look up the closest entry below, which is
  \proc{select$_1$(8)} = 17, in $O(1)$ time. Therefore, we know that the result
  of \proc{select$_1$(14)} must be to the right of \proc{select$_1$(8)}, so we
  scan left to right starting from position 17 in the bitvector until we have
  found 6 more bits (for 8 + 6 = 14 bits total).}

\begin{figure}[t]
    \centering
    \includegraphics[width=.9\columnwidth,page=4,trim=7cm 6cm 9cm 6cm,clip]{visualization.pdf}
    \caption{Examples of rank/select operations using the associated indices on
      top of a bitvector.}
    \label{fig:traditional-bv}
\end{figure}

\rev{\para{Succinct trie \bv structure} BP, LOUDS, and DFUDS are three popular
  \bv-based succinct tree-encoding schemes \cite{jacobson1989space,
    benoit2005representing}. They all use the topology bitvector to map each
  tree element (\rev{i.e.,} nodes, edges, internal nodes, leaf nodes) to a
  unique integer ID in the range of $[0, n-1]$, where $n$ is the number of such
  elements in the tree. Depending on the encoding scheme, the topology may
  support efficient tree operations \rev{(e.g., parent/child navigation)}
  \rev{via \bv operations (e.g., rank/select).} Succinct trees encoded using
  these schemes can be converted to tries by storing data associated with each
  trie edge in a separate array, which is indexed using the edge IDs~\cite{surf,
    marisa, pdt, coco_1}.

  Efficient bitvector operations require auxiliary data structures which
  separate the bitvector into two parts: the \defn{index}, which accelerates
  bitvector operations, and the \defn{bit sequence}, which stores the content of
  the bitvector (i.e. the trie topology). The index facilitates
  \rankopnospace/\selectop operations essential to succinct trie navigation.}

\subsection{LOUDS Encoding}\label{sec:louds}

Level-Order Unary Degree Sequence (LOUDS) is a% a simple yet powerful
\newmymarginpar{R2 O1/D4}\newrev{classical example of} a \bv-based succinct-tree
encoding scheme~\cite{jacobson1989space}. Each tree node is encoded using the
bit sequence $1^{m}0$ (\rev{i.e.}, $m$ 1's followed by a 0), where $m$ is the
degree (\rev{i.e.}, the number of children) of the node. The encoding bits of
each node are concatenated in level order to form the encoding \bv \var{bv} of
the entire tree. With this encoding, each 1 bit corresponds to a parent-to-child
edge, whereas each 0 bit marks the end of a node. We say a position $i$ belongs
to a node $k$ if bit $i$ is part of node $k$'s encoding. The LOUDS encoding
efficiently supports the following tree operations:
\begin{itemize}
    \item \proc{NodeID(i) = bv.rank$_0$(i)}, the LOUDS ID of the node to which position $i$ belongs.
    \item \proc{EdgeID(i) = bv.rank$_1$(i)}, the LOUDS ID of the edge at position $i$ (Requirement: \proc{bv[i] == 1}).
    \item \proc{LeafID(i) = \rev{bv}\mymarginpar{Meta 5, R2 D17}.rank$_{00}$(i)}, the LOUDS ID of the leaf node to which position $i$ belongs.
    \item \proc{InnerID(i) = NodeID(i) - LeafID(i)}, the LOUDS ID of the inner node to which position $i$ belongs.
    \item \proc{Child(i)} = \proc{bv.select$_0$(bv.rank$_1$(i+1)) + 1}, the position of the child node connected to the edge at position $i$.
    \item \proc{Parent(i) = bv.select$_1$(bv.rank$_0$(i))}, the position of the parent of the node to which position $i$ belongs.
\end{itemize}

As illustrated in~\figref{louds-example}, encoding a trie using LOUDS involves
the following steps: First, the \emph{topology} of the trie is encoded using the
LOUDS \bv; Second, all the trie labels are concatenated in level order into a
separate array. At query time, trie labels are uniquely indexed by edge IDs
(because they are associated with trie edges), whereas keys are indexed by leaf
IDs.

\mymarginpar{Meta 1, R1 D1}\rev{Next, we will present an example of how to
  traverse a LOUDS-encoded trie using the trie
  in~\figref{louds-example}. Suppose we are navigating to the second child of
  the root (i.e., from node 0 to node 2). We need to follow the second edge of
  the root node, which is at position 1. From the aforementioned formula,
  \proc{Child(1)} = \proc{bv.select$_0$(bv.rank$_1$(2))+1} =
  \proc{bv.select$_0$(2)+1} = 6. Indeed, if we look at the LOUDS bitvector
  in~\figref{louds-example}, the second child of the root starts at position 6.}

Balanced Parentheses (BP)~\cite{jacobson1989space} and Depth-First Unary Degree
Sequence (DFUDS)~\cite{benoit2005representing} are two popular alternatives to
LOUDS with the same two-bits-per-node overhead, but they traverse the tree in
different orders. Compared with LOUDS, they \rev{support}\mymarginpar{Meta 5, R1
  D9} a broader range of tree operations, including computing subtree sizes or
counting the number of leaf nodes to the left/right of a specific node. However,
\rev{their functionality} comes at the cost of slower parent/child navigation
\rev{compared to LOUDS due to their greater} complexity.

\begin{figure}[t]
    \centering
    \includegraphics[width=.48\textwidth,page=1,trim=7cm 5cm 7cm 4.5cm,clip]{visualization.pdf}

    \caption{LOUDS/LOUDS-Sparse encoding example. Stored keys: car, cat, suc, succ, sum, tie, tip, trie, try.}
    \label{fig:louds-example}
  \end{figure}

\subsection{State-of-the-art Succinct Tries}\label{sec:existing-succinct-tries}

\para{Fast Succinct Trie} The Fast Succinct Trie (FST)~\cite{surf} is a static
succinct trie that combines the benefits of LOUDS-Sparse and LOUDS-Dense by
leveraging both formats in different levels of the trie. Specifically, FST
encodes the top levels of the trie, where nodes tend to have more branches and
are more likely to be queried, using LOUDS-Dense, a representation that is fast
and space efficient for nodes with dense branches. The rest of the trie, which
tends to be colder and of lower average degree, is encoded with the more compact
LOUDS-Sparse. \new{Due to space limitations, we omit the details of
  LOUDS-Dense, %because LOUDS-Sparse is more similar to the schemes that we will focus on,
  but we refer the reader to the original FST paper~\cite{surf} for the full
  details.}

LOUDS-Sparse uses 2 bits to encode each trie edge. The \\\proc{HasChild} bit of
each edge indicates if this edge points to an internal node (1) or a leaf node
(0). The \proc{Louds} bit of each edge indicates if this edge points to the
first child of current node (1) or not (0), which is used to determine node
boundaries. These two sets of encoding bits are respectively concatenated in
level order to form two \bvs, \proc{HasChild} and \proc{Louds}. With
LOUDS-Sparse, we have
\begin{itemize}
    \item \proc{Child(i) = Louds.select$_1$(HasChild.rank$_1$(i+1)+1)}.
    \item \proc{Parent(i) = HasChild.select$_1$(Louds.rank$_1$(i+1)-1)}.
    \item \proc{LeafID(i) = HasChild.rank$_0$(i)}.
\end{itemize}

\rev{Let us consider how to perform trie navigation in a LOUDS-sparse-encoded
  trie using the example in~\figref{louds-example}. Suppose we are navigating
  again to the second child of the root (also at position 1). \mymarginpar{Meta
    1, R1 W1, R1 D1}From above, \proc{Child(1) =
    Louds.select$_1$(HasChild.rank$_1$(2)\\+1) = Louds.select$_1$(3)} = 4. Indeed,
  the second child of the root starts at position 4 in the \proc{Louds}
  bitvector in the LOUDS-Sparse encoding.}

\begin{figure}[t]
    \centering
    \includegraphics[width=\columnwidth,page=2,trim=7cm 5cm 9cm 3.5cm,clip]{visualization.pdf}
    \caption{An example of a CoCo-trie obtained by collapsing sub-tries
      surrounded by dashed boxes in~\figref{louds-example}.}
    \label{fig:coco-example}
\end{figure}

\para{CoCo-Trie} The CoCo-trie \cite{coco_0, coco_1} is a recent
state-of-the-art static succinct trie that achieves both low memory consumption
and fast query performance. A CoCo-trie is constructed by first creating a
regular trie $T$ and then adaptively selecting sub-tries of $T$ to collapse into
macro-nodes. Keys in each collapsed macro-node are transformed to an increasing
sequence of integer codes which is then encoded using a select pool of succinct
integer encoding algorithms. The optimal set of sub-tries to collapse is
computed using a \rev{bottom-up} dynamic-programming
algorithm. % through a bottom-up traversal of the trie. % keeping track of the optimal sub-trie rooted at each node.

\mymarginpar{Meta 1, R1 W1, R1 D1}\rev{~\figref{coco-example} illustrates the
  CoCo-trie encoding of the example trie in~\figref{louds-example}. Since the
  CoCo-trie uses LOUDS encoding, its trie navigation is done through
  \rankop/\selectop operations on the bitvector as described in~\secref{louds}.}


\para{Marisa Trie}\label{sec:marisa} The Marisa trie~\cite{marisa} is a
LOUDS-encoded Patricia trie~\cite{partricia} that contracts unary paths into
single concatenated labels. Unlike Patricia, Marisa improves space efficiency by
recursively storing unary paths in multiple tries.

\new{The upper subgraph of \figref{marisa-example} shows an example of Patricia,
  where unary paths are contracted into a single node. To convert the trie to
  LOUDS succinct form, two bitvectors are needed: the topology bitvector
  \texttt{LOUDS} (\secref{louds}), and an additional bitvector (\texttt{IsLink})
  which indicates if each edge is a regular edge (0) or a contracted edge
  (1). Each contracted edge is uniquely indexed by a \var{LinkID} determined by
  a \rankop operation on the \var{IsLink} \bv. The \texttt{Links} vector stores
  \defn{links}, or references, to the concatenated labels on the contracted
  edge. For example, bit 3 in \texttt{LOUDS} corresponds to bit 2 in
  \texttt{IsLink} and link 0 in \texttt{Links} (``omp'').}

\mymarginpar{R1 D7}\new{Marisa improves on Patricia by recursively storing unary
  paths in secondary tries.} The recursion \rev{continues} until the number of
tries reaches a \rev{preset limit}, at which point the outstanding unary paths
are sorted, deduplicated and moved to a \rev{simple ``sorted'' tail container
  (\secref{text-compression}).} \new{The lower subgraph of
  \figref{marisa-example} shows a worked example of a Marisa trie with one
  recursion. For readability, links between tries are expressed as leaf node IDs
  in this example. Because the second trie only needs to support random access
  but not lookup operations, the keys are stored in reversed form so that they
  can be efficiently retrieved by reading the labels on a bottom-up path. For
  example, to retrieve link 4's corresponding key, we start from leaf 4 of the
  second trie, concatenate all the (reversed) labels on the leaf-to-root path
  (`o' and `mp'), and obtain `omp'. In practice, links are implemented as leaf
  node positions to avoid an additional \selectop operation. Finally, to
  mitigate the cost of link tracing during lookups,
  %because lookup operations at each node may potentially traverse all edges of that node, %which is very expensive because of link tracing,
  Marisa uses a small piece of memory to cache frequently-taken paths.} % to speed up search.}
% A typical Marisa trie first reduces key length through several levels of
% recursions and then stores the outstanding unary paths in tail container.

\begin{figure}[t]
    \centering
    \includegraphics[width=\columnwidth,page=11,trim=8.2cm 5.4cm 10cm 3.9cm,clip]{visualization.pdf}
    \caption{
    % Marisa example. Keys stored: candidate, camera, tie, tip, traffic, trie. DICT is either another Marisa trie or a tail container.\todo{update to illustrate recursion}
      \rev{Examples of Patricia (top) and Marisa (bottom) with one
        recursion. Keys stored: cache, camp, compare, compute. Unary paths in
        the first Marisa trie are stored in the second Marisa trie in reversed
        form: ehc, era, etu, pm, pmo. Links in the first Marisa trie are
        indicated using dashed arrows, and the numbers alongside them indicate
        their values (which in this example are leaf node IDs in the second
        Marisa trie).}}
    \label{fig:marisa-example}
\end{figure}

\new{The number of recursions in Marisa exposes a \emph{space-time tradeoff}, as
  more recursions further compress the data at the cost of degraded search
  performance and higher build times~\cite{marisa-code}.}

\para{Path Decomposed Trie}
The Path Decomposed Trie (PDT) \cite{pdt, pdt_theory, dynamic_pdt} is a
DFUDS-encoded succinct trie. A PDT is obtained by transforming a regular trie
through path decomposition \cite{path_decomp}: At each step, a root-to-leaf path
is selected and contracted into a single node, and the subtrees dangling on the
original path are converted to its child nodes. The procedure \rev{recurses} %is then performed
 on the dangling subtrees.

 The PDT achieves both significant memory reduction and competitive query
 performance by compressing the contracted paths with an approximate version of
 the \defn{\repair} \mymarginpar{Meta 1, R1 W1}\rev{text compression
   algorithm~\cite{larsson2000off, repair}, which we will summarize in the next
   subsection.}

\new{\subsection{Text Compression
     Algorithms}\label{sec:text-compression} \mymarginpar{R1
     D7}\figref{repair-example} presents a worked example of how to compress a
   set of strings with the sorted tail container, \repair algorithm, and
   \fsst. Next, we summarize these algorithms.

  \para{Sorted tail container} \mymarginpar{Meta 1}The sorted tail container is
  the string container for the last Marisa trie on a recursion chain. \rev{It
    overlaps two keys if one is a suffix of the other, and is particularly space
    efficient for datasets with short and repetitive keys. To efficiently detect
    possibilities of overlapping, it needs to reverse and sort all keys. Once
    the strings have been sorted, construction takes linear time.} \new{The link
    array which is used to retrieve keys from the container should also preserve
    the original key ordering.}

  \new{\para{Dictionary compression} Dictionary compression encodes common text patterns as small integer
    codes, exemplified by the famous LZW algorithm~\cite{lzw}. The mapping
    between patterns and integers is known as the \defn{dictionary}, and finding
    a good dictionary is known to be a key challenge to such
    algorithms~\cite{fsst}. In our work, we focus on two related algorithms that
    support fast random access: re-pair~\cite{larsson2000off} and
    FSST~\cite{fsst}.}

  \para{Exact and approximate \repair} At a high level, the original \repair
  algorithm~\cite{larsson2000off} replaces the most frequently-occurring
  character pair with a new special character, proceeding in rounds. \rev{For
    example, using special characters $\alpha$, $\beta$, $\gamma$ and $\delta$,
    \repair may compress the string aabbaabb as $\alpha$bb$\alpha$bb ($\alpha$ =
    aa) in the first pass, $\alpha\beta\alpha\beta$ ($\beta$ = bb) in the second
    pass, $\gamma\gamma$ ($\gamma = \alpha\beta$) in the third pass and $\delta$
    ($\delta = \gamma\gamma$) in the fourth pass.} \new{The dictionary is
    expressed as a set of recursive rules of pairing, which can be flattened to
    support constant-time decoding~\cite{pdt}, with a small space
    overhead. \mymarginpar{R1 D4}Prior work shows that \repair compresses a
    sequence $T$ of length $n$ over an alphabet of size $\sigma$ and $k$-th
    order entropy $H_k$ to $O(H_k)$ bits.~\cite{navarro2008re,
      ochoa2018repair}.}

\begin{figure}[t]
    \centering
   % \begin{minipage}[h]{.48\textwidth}
        \includegraphics[width=\columnwidth,page=14,trim=7.5cm 11.5cm 11.5cm 2cm,clip]{visualization.pdf}
        \caption{
        \rev{Examples of re-pair, FSST and sorted tail containers. Re-pair and FSST are both dictionary encoding methods (denoted with \texttt{Dict}) that replaces common text patterns with small integer codes. The sorted tail container (denoted with \texttt{Sorted}) overlaps suffix keys.}}
        \label{fig:repair-example}
  \end{figure}

  % \para{Approximate \repair}

  In practice, the PDT uses an \emph{approximate} version of the \repair
   algorithm~\cite{repair} that accelerates build time by orders of magnitude compared to exact \repair by identifying and
   replacing the top $k$ %(for some $k$)
   most frequently-seen character pairs in
   each round, rather than one at a time. To our knowledge, there is no
   theoretical bound on how far the compression ratio of approximate
   \repair is from exact \repair's, as it depends on the dataset.


   \para{Fast Static Symbol Table} The Fast Static
   Symbol \\Table (\fsst)~\cite{fsst} is similar to approximate \repair, but
   optimizes for \emph{lightweight} compression (with linear build
   times). Specifically, it restricts the size of the symbol table (i.e., the
   list of special characters for replacement) to 256 entries and expands the
   size of each replaced entry to substrings of 8 bytes (rather than two
   characters in the original approximate \repair). \fsst further optimizes for
   build time by first encoding
   %running the compression algorithm on
   a small sample (about 16KB) of the original dataset, and then using the
   resulting symbol table to encode the entire dataset. FSST implements
   hardware-specific optimizations, achieving extremely high compression and
   decompression throughput. In practice, \fsst reduces the compression time by
   over an order of magnitude compared to approximate \repair, while achieving
   similar compression ratios. However, \newmymarginpar{R2 O1/D4}\newrev{no
     theoretical space bound has been proved for FSST}.}
%%% Local Variables:
%%% mode: latex
%%% TeX-master: "main"
%%% End:
